\begin{center}
	\zihao{4}\heiti
	% 中国科学院大学\\
	\LARGE{中国科学院研究生院\\$\boldsymbol{2009}$年招收硕士研究生入学统一考试试卷\\
	科目名称：电动力学}

\end{center}
\noindent
考生须知：\\
1.本试卷满分为150分,全部考试时间总计180分钟。\\
2.所有答案必须写在答题纸上,写在试题纸上或草稿纸上一律无效。

\vspace*{2ex}
\hrule
\vspace*{4ex}

\fancypagestyle{mypagestyle}{%
	\fancyhf{}% Clear header/footer
	\fancyfoot[OR]{第~\thepage~页，共2页}%
	\fancyfoot[EL]{第~\thepage~页，共2页}%
	\fancyfoot[OL]{科目名称：电动力学 }%
	\fancyfoot[ER]{科目名称：电动力学 }%
	\renewcommand{\headrulewidth}{0pt}% Header rule of .4pt
	\renewcommand{\footrulewidth}{0.5pt}% Header rule of .4pt
}
\fancypagestyle{plain}{%
	\fancyhf{}% Clear header/footer
	\fancyfoot[OR]{第~\thepage~页，共2页}%
	\fancyfoot[EL]{第~\thepage~页，共2页}%
	\fancyfoot[OL]{科目名称：电动力学 }%
	\fancyfoot[ER]{科目名称：电动力学 }%
}

\setcounter{page}{1}%设置页面计数器
\pagestyle{mypagestyle}

\onehalfspacing
% 排版结束，下面是题目

\vspace{0.4cm}

\noindent
一、填空和选择题（40分）
\vspace{0.2cm}

\noindent
1.\hspace{0.8em}真空中电磁波的能流密度矢量$\Vec{S}$和动能密度矢量$\Vec{ \mathrm{g} }$的表达式为$\Vec{S}=\underline{  \qquad  \qquad  \qquad  \qquad \qquad }$，
\vspace{0.3cm}

\noindent
$\Vec{ \mathrm{g} }=\underline{  \qquad  \qquad  \qquad  \qquad \qquad }$.
\vspace{0.3cm}

\noindent
2.\hspace{0.8em}海水的$\mu_{\gamma}=1$，$\sigma=1$(欧姆·米)$^{-1}$，($\mu_0 = 4 \pi \times 10^{-7}$享/米)，则频率$\nu$为$10^6$赫兹的电磁波，
\vspace{0.2cm}

\noindent
在海水中的透入深度为\underline{  \qquad  \qquad  \qquad  \qquad \qquad }米
\vspace{0.2cm}

\noindent
3.\hspace{0.8em}在$\displaystyle{ \frac{\sigma}{\omega \varepsilon} } \gg 1 $极限情况下，透入金属的电磁波的磁场比电场矢量的位相要落后，其位相差$\varphi=$ \vspace{0.2cm}

\noindent
\underline{  \qquad \qquad \qquad \qquad \qquad \qquad}

\noindent
4.\hspace{0.8em}试用相对论四维形式写出麦克斯韦方程组\underline{  \qquad \qquad \qquad  \qquad \qquad \qquad \qquad  \qquad \qquad \qquad \qquad  }

\vspace{0.3cm}
\noindent
和电荷守恒定律\underline{  \qquad \qquad \qquad \qquad \qquad \qquad  \qquad }

\vspace{0.3cm}
\noindent
5.\hspace{0.8em}电磁波在无界真空和波导管中传播时

\vspace{0.2cm}

A.\hspace{0.8em}均有$v_{\text{群速}} \cdot v_{\text{相速}} = c^2$ \hspace{5em}
B.\hspace{0.8em}均有$v_{\text{相速}} = c$  \hspace{3em}
\vspace{0.1cm}

C.\hspace{0.8em}均为TEM波 \hspace{7.9em}
D.\hspace{0.8em}均有截止频率 \vspace{0.2cm}

\noindent
6.\hspace{0.8em}在均匀介质内部的体极化电荷密度$\rho_P$总是等于自由电荷密度$\rho_f$的
\vspace{0.2cm}

A.\hspace{0.8em} $-(1 - \displaystyle{  \frac{1}{ \varepsilon_r } } )$倍 \hspace{8em}
B.\hspace{0.8em} $ (1 - \displaystyle{  \frac{1}{ \varepsilon_r \varepsilon_0 } } )$倍
\vspace{0.2cm}

C.\hspace{0.8em} $ (1 - \displaystyle{  \frac{1}{ \varepsilon_0 } } )$倍 \hspace{8.9em}
D.\hspace{0.8em} $ -(1 - \displaystyle{  \frac{1}{ \varepsilon_0 } } )$倍

\vspace{0.2cm}

\noindent
7.\hspace{0.8em}均匀带电的长形旋转椭球体带总电量Q，坐标原点在椭球中心，其
\vspace{0.2cm}

A.\hspace{0.8em}电偶极矩和电四极矩与$Q^2$有关 \hspace{2em} B.\hspace{0.8em}电偶极矩为零，电四极矩不为零
\vspace{0.2cm}

C.\hspace{0.8em}电偶极矩不为零，电四极矩为零\hspace{2.1em} D.\hspace{0.8em}电偶极矩，电四极矩都为零
\vspace{0.2cm}

\noindent
8.\hspace{0.8em}在两个夹角为$60^{\circ}$的接地导体平板内有一点电荷q，那么用电象法求解此区间的电场时，其象电荷的数目为
\vspace{0.2cm}

A.\hspace{0.8em} 三个\hspace{5em}B.\hspace{0.8em}四个\hspace{5em}C.\hspace{0.8em} 五个 \hspace{5em}D.\hspace{0.8em} 六个
\vspace{0.2cm}

\noindent
9.\hspace{0.8em}如图，一个单位面积的法线矢量$d \Vec{\sigma}$与$\Vec{E}$夹角为$45^{\circ}$，当$\Vec{B} = 0$，$\Vec{E} \neq 0$时作用在该面元$d \Vec{\sigma}$ \vspace{0.1cm}

\noindent
上的电场力$d \Vec{f_{E}}$ 与$d \Vec{\sigma}$ 的夹角为：（顺时针为正，反时针为负）



\vspace{0.5cm}

A.\hspace{0.8em} $0^{\circ}$\hspace{8em}  B.\hspace{0.8em} $22.5^{\circ}$
\vspace{0.3cm}

C.\hspace{0.8em} $45^{\circ}$\hspace{7.5em}  D.\hspace{0.8em} $90^{\circ}$ 

\vspace{-2.8cm}
\begin{figure}[h!]
    \flushright
    \includegraphics[height=2.8cm]{test/2009/pic_1.png}
\end{figure}

\newpage
\noindent
10.\hspace{0.8em}一电磁波垂直入射到一个理想的导体平面上，则有
\vspace{0.1cm}

A.\hspace{0.8em} 反射波的$\Vec{E}$矢量的位相要改变$180^{\circ}$\vspace{0.1cm}

B.\hspace{0.8em} 反射波的$\Vec{H}$矢量的位相要改变$180^{\circ}$\vspace{0.1cm}

C.\hspace{0.8em} 反射波的$\Vec{E}$矢量和$\Vec{H}$矢量的位相都要改变$180^{\circ}$\vspace{0.1cm}

D.\hspace{0.8em} 反射波的$\Vec{E}$矢量和$\Vec{H}$矢量的位相都不改变\vspace{0.5cm}

\noindent
二、\hspace{0.2em}在无限大各向同性均匀介质中，挖一半径为$R_0$的球状真空，介质的介电系数为$\epsilon$，把它放在一个均匀外电场$\Vec{E_0}$中，求球内外的电势。（20分）
\vspace{3cm}

\noindent
三、\hspace{0.2em}XY平面上有一均匀带电、半径为a的圆环，中心在原点处，其总电荷量为+Q，求此系统的电偶极矩、电四极矩及在远处的电势$\Phi$（精确到四极矩项）。（20分）
\vspace{3cm}

\noindent
四、\hspace{0.2em}带电粒子（电荷为q）在XY平面上绕Z轴作匀速圆周运动，角频率为$\omega$，半径为$R_0$，设$\omega R_0 \ll c$，试计算辐射场公布与能量流密度。（20分）

\vspace{3cm}

\noindent
五、\hspace{0.2em}在坐标系$\Sigma$中，有两个物体都以速度u沿X轴运动，在$\Sigma$系看来，它们一直保持距离L不变。今有一观察者以速度v沿X轴运动，他看到此二物体的距离是多少？\vspace{3cm}

\noindent
六、\hspace{0.2em}超导体中的电流密度$\Vec{J_s}(r)$与其中的电磁场矢量势$\Vec{A}(r)$的关系为$\Vec{J_s} = - \displaystyle{ \frac{1}{c \Lambda} } \Vec{A}$，其中$\Lambda$为常\vspace{0.1cm}

\noindent
数，若超导体的磁导率为$\mu$，且内部电场不随时间变化，利用麦克斯韦方程导出超导体中磁感应强度所满足的方程（不包含$\Vec{J_s}$）。在一维半无限超导体的情况下，若在端点的磁感应强度为$\Vec{B}(0)$，试求磁感应强度在此一维半无限超导体中的表达式。（30分）